\section*{第 \ref{ch:g7:triangles} 章 --- 三角形与角}

\begin{solution}{exo:g7:triangles:1}
余角：$60^\circ$；$45^\circ$；$18^\circ$。
补角：$150^\circ$；$135^\circ$；$108^\circ$。
\end{solution}

\begin{solution}{exo:g7:triangles:2}
对顶角的度数也是 $35^\circ$
（见\cref{thm:g7:triangles:equalities}）；其余两个角均与它互补，
各为 $180 - 35 = 145^\circ$。
\end{solution}

\begin{solution}{exo:g7:triangles:3}
(a) $180 - 100 = 80^\circ$. \quad
(b) $180 - 118 = 62^\circ$. \quad
(c) $180 - 150 = 30^\circ$.
\end{solution}

\begin{solution}{exo:g7:triangles:4}
顶角为 $40^\circ$：两底角之和为 $180 - 40 = 140^\circ$，
二者相等，各为 $70^\circ$。

底角为 $70^\circ$：两底角相等，故顶角为
$180 - 2 \times 70 = 40^\circ$。
\end{solution}

\begin{solution}{exo:g7:triangles:5}
$(6, 8, 12)$：$12 < 6 + 8 = 14$，能构成三角形。

$(5, 5, 11)$：$11 > 5 + 5 = 10$，不能构成三角形。

$(4, 9, 13)$：$13 = 4 + 9$，是“三角形”的退化情形 --- 三点共线，
不是真正的三角形。

$(7, 7, 7)$：$7 < 14$，是等边三角形。
\end{solution}

\begin{solution}{exo:g7:triangles:6}
$x + 2x + 3x = 6x = 180^\circ$，故 $x = 30^\circ$，三个角分别为
$30^\circ$、$60^\circ$ 和 $90^\circ$ --- 是直角三角形。
\end{solution}

\begin{solution}{exo:g7:triangles:7}
预测：$\widehat{BAC} = 180 - (50 + 60) = 70^\circ$；
完成作图后用量角器测量可验证此结果。
\end{solution}

\begin{solution}{exo:g7:triangles:8}
在第一个交点处，四个角依次为 $54^\circ$、$126^\circ$、
$54^\circ$、$126^\circ$（对顶角与邻补角）。
平行线在第二个交点处复现相同的四个度数
（同位角）：共八个角，其中四个为 $54^\circ$，
四个为 $126^\circ$。
\end{solution}

\begin{solution}{exo:g7:triangles:9}
设 $c$ 为第三边。三角形不等式要求
$c < 8 + 3 = 11$，且 $8 < c + 3$，即 $c > 5$。于是
$5 < c < 11$，故长为 $7$ cm 的边可行；长为 $12$ cm（或 $4$ cm）的
边则不行。
\end{solution}

\begin{solution}{exo:g7:triangles:10}
$\widehat A + \widehat B = 180 - 90 = 90^\circ$，且
$\widehat A = 2\widehat B$，于是 $3\widehat B = 90^\circ$：
则 $\widehat B = 30^\circ$，$\widehat A = 60^\circ$。
\end{solution}

\begin{solution}{exo:g7:triangles:11}
对角线 $[AC]$ 将 $ABCD$ 分成三角形 $ABC$ 与 $ACD$。
四边形的四个角恰是两个三角形的六个角，
只是重新组合（$A$ 与 $C$ 处的角各被一分为二）。
合计 $180 + 180 = 360^\circ$。五边形从一个顶点出发可分成三个三角形，
共 $3 \times 180 = 540^\circ$。
\end{solution}

\begin{solution}{pb:g7:triangles:1}
\textbf{1.} 从六边形的一个顶点引出对角线，可将其分成
$4$ 个三角形，故内角和为
$4 \times 180 = 720^\circ$。

\textbf{2.} 从一个顶点出发，对角线连向除自身及两个相邻顶点外的所有顶点，
共 $n - 3$ 条，将多边形分成
$n - 2$ 个三角形。多边形的每个内角都分配给这些三角形，
故其内角和为
\[
(n - 2) \times 180^\circ .
\]
验证：$n = 3$ 时为 $180^\circ$；$n = 4$ 时为 $360^\circ$
（见\cref{exo:g7:triangles:11}）；$n = 5$ 时为 $540^\circ$；$n = 6$ 时
为 $720^\circ$。

\textbf{3.} 十边形：$8 \times 180 = 1\,440^\circ$。十二边形：
$10 \times 180 = 1\,800^\circ$。

\textbf{4.} 除以 $n$：三角形 $60^\circ$；正方形
$90^\circ$；五边形 $\frac{540}{5} = 108^\circ$；六边形
$\frac{720}{6} = 120^\circ$；八边形 $\frac{1080}{8} =
135^\circ$.

\textbf{5.} 增加一个顶点，扇形剖分中便多一个三角形：
新多边形需要 $n - 1$ 个三角形，而原来需要
$n - 2$ 个。多出一个三角形，内角和便多 $180^\circ$。

\textbf{6.} 行走者转过的角等于内角的补角：
$180^\circ - \widehat A$。对于 $40$--$60$--$80$ 的三角形，
转角依次为 $140^\circ$、$120^\circ$ 和 $100^\circ$。

\textbf{7.} $140 + 120 + 100 = 360^\circ$：恰好转过一整圈。

\textbf{8.} 绕行一周之后，你的鼻子起始与终止时指向同一方向，
恰好转过一整圈：所有拐角处转过的角之和
为一个周角，
$360^\circ$ --- 无论拐角有多少个，街区是什么形状。
（这一事实与 $n$ 无关，十分奇妙。）

\textbf{9.} 每个拐角处，内角 $+$ 转角 $= 180^\circ$。
对 $n$ 个拐角求和：
\[
\text{（内角之和）} + 360^\circ = 180^\circ \times n,
\qquad\text{故}\qquad
\text{内角之和} = (n - 2) \times 180^\circ
\]
--- 即第 2 题的公式，用行走方法再次证明。

\textbf{10.} 内角为 $150^\circ$，说明每个转角为
$30^\circ$，于是 $n = \frac{360}{30} = 12$：正十二边形。
若内角为 $155^\circ$，则转角为 $25^\circ$，
此时 $\frac{360}{25} = 14.4$ 不是整数个拐角，
故不存在这样的正多边形。

\textbf{11.} 三角形（$60^\circ$）：$6$ 个相接，
$6 \times 60 = 360$。正方形（$90^\circ$）：$4$ 个相接。六边形
（$120^\circ$）：$3$ 个相接，$3 \times 120 = 360$。三者
都恰好严丝合缝。

\textbf{12.} 三个正五边形相接：$3 \times 108 = 324^\circ$
--- 尚余 $36^\circ$ 的空隙。四个相接：
$4 \times 108 = 432^\circ > 360^\circ$ --- 发生重叠。两种情形
均不可行：正五边形不能铺满平面。

\textbf{13.} 在铺砌的某个拐角处，至少有 $3$ 块图形相接。
边数多于六的正多边形，其内角 \emph{大于} $120^\circ$
（第 4 题及内角随 $n$ 增大的规律），
因此三个拐角处的角之和超过
$3 \times 120 = 360^\circ$：发生重叠，不可能。恰为六边时，
$3 \times 120 = 360$ 恰好可行；少于六边时，由第 11 题与第 12 题分类讨论。
正则镶嵌的完整列表为：
三角形、正方形、六边形。

\textbf{14.} 八边形与正方形：
$135 + 135 + 90 = 360^\circ$：拐角恰好闭合 --- 这是经典的
街道铺法。正方形、六边形、正十二边形：正十二边形的
转角为 $\frac{360}{12} = 30^\circ$，故其内角为
$150^\circ$，且 $90 + 120 + 150 = 360^\circ$：同样可行。

\textbf{15.} 蜂巢的拐角：$3 \times 120 = 360^\circ$，
恰好无缝。在面积相等的三角形、正方形与六边形之中，
六边形的周长最短 --- 在三者中它
最接近狄多的圆，
（\cref{pb:g6:measure:1}） --- 故六边形巢室以最少的蜡围住同样多的蜂蜜：
蜜蜂像几何学家一样铺砌。
\end{solution}
